\section{Restriction and the prescribed data at a special fiber}
\label{data:section}

Let $(S,h)$ now be any fiber under consideration, and put
\[
 V=h^\perp\subset H^2(S,\mathbb Q),\qquad
 T=\operatorname{NS}(S)_{\mathbb Q}^{\perp}\subset V.
\]
Suppose that the primitive Clifford multiplication for the cup form
$q$ on $V$ has been algebraically realized. We show that this gives
the stipulated standard map on $T$, including its chosen rational
quadratic polarization. The arguments use no hypothesis on the
Hodge group at this fiber.

\begin{proposition}[Restriction to the transcendental part]
\label{data:restriction}
Set $q_0=q|_T$. Algebraic realization of $m^q$ on $V$ implies
algebraic realization of $m^{q_0}$ on $T$.
\end{proposition}

\begin{proof}
The orthogonal decomposition $V=T\oplus T^\perp$ is a decomposition
of rational Hodge structures, and $T$ contains the period plane.
By the Clifford Basis Theorem there are inclusions
\[
 C^\pm(T,q_0)\hookrightarrow C^\pm(V,q).
\]
They are weight-one Hodge inclusions: the element $J$ defining the
left Hodge circle already belongs to $C^+(T_{\mathbb R},q_0)$.
Let $i:C^+(T,q_0)\to C^+(V,q)$ be the inclusion and choose a rational
Hodge retraction $p:C^-(V,q)\to C^-(T,q_0)$. Such a retraction
exists by polarizability: the orthogonal complement for a
weight-one polarization is a rational Hodge complement.
For every $x\in T$,
\begin{equation}\label{data:restriction-identity}
 p\,L_x^{V}\,i=L_x^{T},
\end{equation}
because $xC^+(T,q_0)\subset C^-(T,q_0)$ and $p$ restricts to the
identity there. Lemma~\ref{cliff:slots} realizes the two operations
on the abelian factors. Restricting the resulting cycle action to
$T$ gives the asserted equality; no operation on the K3 factor
is needed.
\end{proof}

\begin{proposition}[Change of quadratic polarization]
\label{data:polarization}
Let $q_0$ and $q_1$ be rational quadratic Hodge polarizations on
the same K3-type Hodge structure $T$, allowing either overall sign
as in Section~\ref{cliff:section}. If $T\subset H^2(S,\mathbb Q)$
and $m^{q_0}$ is algebraically realized, then so is $m^{q_1}$.
\end{proposition}

\begin{proof}
\emph{Step 1: square-root isometries.}
Write
\begin{equation}\label{data:P}
 q_1(u,v)=q_0(Pu,v).
\end{equation}
The two form isomorphisms from $T$ to its twisted dual show that
$P$ is an invertible rational Hodge endomorphism. Symmetry of the
forms makes $P$ self-adjoint for $q_0$.
Let $K$ be $+1$ on the real period plane and $-1$ on the real
$(1,1)$-part, and let $\epsilon_0$ be the sign of $q_0$ on the
period plane. Then
\[
 (u,v)_0=\epsilon_0 q_0(Ku,v)
\]
is positive definite. Since $P$ commutes with $K$, it is also
self-adjoint for this positive definite real form. Hence $P$ is
semisimple with real nonzero eigenvalues. Its distinct eigenvalues
$\lambda_1,\ldots,\lambda_r$ and spectral projectors
$p_1,\ldots,p_r$ are defined over $\overline{\mathbb Q}$, the latter
being polynomials in $P$.

Choose $\nu_i\in\overline{\mathbb Q}$ with $\nu_i^2=\lambda_i$.
For each $s=(s_1,\ldots,s_r)\in\{1,-1\}^r$, put
\begin{equation}\label{data:roots}
 R_s=\sum_{i=1}^r s_i\nu_i p_i.
\end{equation}
These are invertible self-adjoint operators with $R_s^2=P$.
Consequently
\[
 q_0(R_su,R_sv)=q_0(Pu,v)=q_1(u,v).
\]
Thus $R_s:(T,q_1)_{\overline{\mathbb Q}}\to
(T,q_0)_{\overline{\mathbb Q}}$ is an isometry and extends to a
parity-preserving Clifford algebra isomorphism
\[
 F_s:C(T,q_1)_{\overline{\mathbb Q}}
          \xrightarrow{\sim}C(T,q_0)_{\overline{\mathbb Q}}.
\]
\smallskip\noindent
\emph{Step 2: weight-one Hodge maps.}
We show that the parity restrictions of $F_s$ and their inverses lie
in the $\overline{\mathbb Q}$-span of rational weight-one Hodge maps.

Let $\mathcal J$ be the joint rational Hodge group of $T$ and the
four weight-one spaces $P_j^\pm=C^\pm(T,q_j)$, $j=0,1$.
Here the joint Hodge group is the smallest rational algebraic
subgroup containing their simultaneous real Hodge circle.
Its geometric identity component is defined over $\mathbb Q$ and
already contains that connected circle. Minimality therefore gives
$\mathcal J=\mathcal J^\circ$, so $\mathcal J$ is geometrically
connected. Write $\rho_T$ for its action on
$T$. Since $P$ is a rational Hodge endomorphism, it commutes with
$\rho_T(\mathcal J)$; therefore every $R_s$ does as well after
scalar extension.

For each $j$, the projection of $\mathcal J$ to the actions on
$(T,P_j^+,P_j^-)$ lies in the image of
\[
 \operatorname{Spin}(T,q_j)\longrightarrow
 \operatorname{GL}(T)\times\operatorname{GL}(P_j^+)
                      \times\operatorname{GL}(P_j^-),
\]
using the orthogonal action on $T$ and left multiplication on
the parity spaces. Indeed, this rational algebraic image contains
the projected circle by \eqref{cliff:circle}. The representation
is faithful, since the left action on $C^+$ recovers a spin
element by evaluation at $1$. Thus it identifies the spin group
with its algebraic image, and gives algebraic homomorphisms
\[
 \rho_j:\mathcal J\longrightarrow\operatorname{Spin}(T,q_j)
\]
covering the same action $\rho_T$.

The isometry $R_s$ induces an isomorphism of spin groups
$\widehat R_s:\operatorname{Spin}(T,q_1)_{\overline{\mathbb Q}}
\to\operatorname{Spin}(T,q_0)_{\overline{\mathbb Q}}$.
Since $R_s$ commutes with $\rho_T$, both
$\widehat R_s\rho_1$ and $\rho_0$ lift the same orthogonal
representation of $\mathcal J_{\overline{\mathbb Q}}$.
Their pointwise quotient therefore takes values in the central
kernel $\mu_2$ of the spin covering. This quotient is a morphism
from a connected group to a finite group, and is the identity at
the identity. Hence
\begin{equation}\label{data:lift-agreement}
 \widehat R_s\rho_1=\rho_0.
\end{equation}
Because $F_s$ is an algebra isomorphism,
\eqref{data:lift-agreement} says precisely that $F_s$ intertwines
the left $\mathcal J$-actions on both parities.

For finite-dimensional rational representations, equivariant
Hom-spaces commute with field extension: equivariance is a set
of rational linear equations on the matrix of a map. Rational
maps intertwine $\mathcal J$ exactly when they intertwine its
Hodge circle. Thus
\begin{equation}\label{data:hom-descent}
 \operatorname{Hom}_{\mathcal J}(P_1^\pm,P_0^\pm)
       \otimes_{\mathbb Q}\overline{\mathbb Q}
 =\operatorname{Hom}_{\mathcal J_{\overline{\mathbb Q}}}
       (P_{1,\overline{\mathbb Q}}^\pm,
        P_{0,\overline{\mathbb Q}}^\pm)
\end{equation}
expresses $F_s^\pm$, and likewise $(F_s^\pm)^{-1}$, as finite
linear combinations of rational weight-one Hodge maps.

\smallskip\noindent
\emph{Step 3: recovery and rational descent.}
Fix abelian realizations and rational dualities for the $q_1$
target, and let $\mathcal A$ be the rational subspace of maps
\[
 T\longrightarrow P_1^-\otimes P_1^+
\]
obtained by restricting rational codimension-two cycle actions
from $S$ and projecting to this K\"unneth component. It is a
subspace of a finite-dimensional rational Hom-space.
By Lemma~\ref{cliff:slots} and \eqref{data:hom-descent}, conjugating
the algebraically realized $q_0$ multiplication on its Hom slots
gives elements of $\mathcal A\otimes\overline{\mathbb Q}$.
The exact Clifford identity is
\begin{equation}\label{data:conjugation}
 (F_s^-)^{-1}L_x^{q_0}F_s^+
       =L_{R_s^{-1}x}^{q_1}\qquad(x\in T).
\end{equation}
Here the conjugation is expanded using rational Hodge maps on
the two abelian factors; it does not apply $R_s^{-1}$ to the
cohomology of $S$ by a correspondence.

An elementary sign average recovers the unmodified source:
\begin{equation}\label{data:sign-average}
 p_i=\frac{\nu_i}{2^r}
        \sum_{s\in\{1,-1\}^r}s_iR_s^{-1},
 \qquad \operatorname{id}_T=\sum_{i=1}^r p_i.
\end{equation}
Indeed $R_s^{-1}=\sum_j s_j\nu_j^{-1}p_j$ and
$\sum_s s_i s_j=2^r\delta_{ij}$.
Equations~\eqref{data:conjugation} and \eqref{data:sign-average}
therefore place the tensor realization of $x\mapsto L_x^{q_1}$
in $\mathcal A\otimes\overline{\mathbb Q}$.
This map is rational, so Lemma~\ref{cyc:descent} places it in
$\mathcal A$ itself. A finite rational linear combination of
actual cycles realizes it.
\end{proof}

\begin{remark}\label{data:scope}
The proof never uses an algebraic correspondence representing
$P$, any $p_i$, or any $R_s$ on the K3 surface. They enter only
the identities that produce weight-one intertwiners and the
finite rational descent. In particular $P$ need not be scalar.
The argument applies with additional real multiplication or
complex multiplication endomorphisms, and requires no description of the special
Hodge group. It applies also in ranks $2$, $3$, and $4$.
In rank $2$, the full even-Clifford space is two-dimensional and
$\operatorname{Spin}_2\to\operatorname{SO}_2$, over an algebraic
closure, is the squaring map on $\mathbb G_m$ with kernel
$\mu_2$. The faithful left Clifford action and the connected-lift
argument remain valid. Higher small ranks introduce no exception.
We make no assertion that all Hodge embeddings for a smaller
special Hodge group are generated by these operations.

The opposite-sign case is directly visible in rank two. Take
$q_0=\operatorname{diag}(1,1)$ and $q_1=-q_0$ in the same period
basis. Then $J_1=-e_1e_2$ and the isometry $R=i\operatorname{id}$
induces $F(J_1)=-(ie_1)(ie_2)=J_0$. Thus the nonreal square root
preserves the specified weight-one circle, as the lift argument predicts.
\end{remark}

For completeness, a rational Hodge automorphism $D$ of $T$ in a
prescribed identification can also be incorporated without
applying it as a K3 correspondence. First use
Proposition~\ref{data:polarization} for the pullback form $D^*q_1$.
The rational isometry $D:(T,D^*q_1)\to(T,q_1)$ induces
weight-one Hodge isomorphisms on Clifford spaces: the same
connected-lift argument proves this, since $D$ commutes with the
joint Hodge action. Transport on the two Hom slots then gives
exactly $x\mapsto L_{Dx}^{q_1}$.

Finally apply Proposition~\ref{data:restriction}, then
Proposition~\ref{data:polarization}, and then the prescribed
right multiplier and duality of \eqref{cliff:standard}.
Lemma~\ref{cliff:slots} transports the resulting tensor to the
actual $H^1(A_S,\mathbb Q)$ model. The target K\"unneth projectors
give a cycle
\[
 \Gamma_S\in\operatorname{CH}^2(S\times A_S\times A_S)_{\mathbb Q}
 \quad\text{with}\quad
 (\Gamma_S)_*x=\kappa_S(x)\quad\text{for every }x\in T.
\]
This is equality with the prescribed full map, including its
rational normalization, and completes the exact-data step in
Theorems~\ref{main:generic} and~\ref{main:specialization}.
